\section{序列建模的实际应用：语言建模}

\subsection{Next Word Prediction 任务}

语言建模（Language Modeling）是序列建模最典型、最重要的应用之一。其核心任务极其简单：\textbf{给定一系列单词，预测下一个最可能的单词}。

\begin{figure}[H]
\centering
\includegraphics[width=0.85\textwidth]{slides-images/slide-35.jpg}
\caption{预测下一个词：语言建模的核心任务\protect\footnotemark}
\end{figure}
\footnotetext{Slides 第35页。}

例如给定句子 ``This morning I took my cat for a''，模型需要预测下一个词是 ``walk''。

\begin{importantbox}{Next Word Prediction 是 LLM 的基础}
今天的大语言模型（如 GPT 系列）本质上就是在执行这个 Next Word Prediction 任务。通过在海量文本上训练预测下一个词的能力，模型学会了语言的结构和语义。
\end{importantbox}

\subsection{如何将语言表示为数字}

神经网络只能处理数值数据，因此我们需要将单词转化为数值向量。这个过程分为两步：

\textbf{Step 1: 建立词汇表（Vocabulary）}

定义一个包含所有可能单词的有限词汇表，并为每个单词分配一个唯一的索引编号：
\begin{itemize}
    \item ``a'' $\rightarrow$ 1
    \item ``cat'' $\rightarrow$ 2
    \item ``the'' $\rightarrow$ 3
    \item $\ldots$
\end{itemize}

\textbf{Step 2: Embedding（嵌入）}

将索引映射为固定维度的向量。最简单的方法是 \textbf{One-Hot Encoding}：

$$\text{``cat''} \rightarrow [0, 1, 0, 0, \ldots, 0]$$

\begin{warningbox}{One-Hot Encoding 的局限}
One-Hot Encoding 有两个严重缺陷：
\begin{enumerate}
    \item \textbf{维度灾难}：向量维度等于词汇表大小，对于包含数万甚至数十万个词的词汇表，这种表示极其稀疏和低效
    \item \textbf{缺乏语义信息}：所有单词之间的``距离''相等，``cat'' 和 ``dog'' 的距离与 ``cat'' 和 ``airplane'' 相同，无法捕捉语义相似性
\end{enumerate}
\end{warningbox}

更好的方法是使用\textbf{Learned Embedding}（可学习嵌入）：通过训练让网络自动学习一种低维、稠密的向量表示，使语义相近的词在向量空间中距离较近。

\begin{figure}[H]
\centering
\includegraphics[width=0.75\textwidth]{slides-images/slide-38.jpg}
\caption{Learned Embedding：语义相近的词在向量空间中距离较近\protect\footnotemark}
\end{figure}
\footnotetext{Slides 第38页。}

\begin{knowledgebox}{Embedding 的本质}
Embedding 本质上是一个可学习的查找表（Lookup Table）。它将离散的符号（如单词、字符）映射到连续的向量空间中。经过训练后，语义相似的词（如 ``king'' 和 ``queen''）在嵌入空间中会彼此接近，而语义不同的词则相距较远。这种特性对于下游的序列建模任务至关重要。
\end{knowledgebox}

\subsection{序列中的短期与长期依赖}

自然语言中存在不同范围的依赖关系：

\begin{itemize}
    \item \textbf{短期依赖}：``The clouds are in the \underline{\hspace{1cm}}'' $\rightarrow$ ``sky''。答案仅依赖前面几个词。
    \item \textbf{长期依赖}：``I grew up in France, $\ldots$ and I speak fluent \underline{\hspace{1cm}}'' $\rightarrow$ ``French''。答案依赖很久之前的上下文``France''。
\end{itemize}

长期依赖对 RNN 提出了严峻挑战，这直接引出了下一节的主题。

\subsection{本章小结}

语言建模的核心任务——预测下一个词——看似简单却极其强大，它是所有现代大语言模型的训练基础。将语言表示为数值向量（Embedding）是实现这一任务的第一步，而处理序列中的长短期依赖关系则是核心挑战。

\newpage
